\documentclass[journal]{IEEEtran}

\usepackage{amsmath,amssymb,amsfonts}
\usepackage{graphicx}
\usepackage{booktabs}
\usepackage{cite}
\usepackage{microtype}
\usepackage{url}

\setlength{\textfloatsep}{6pt plus 1pt minus 1pt}
\setlength{\floatsep}{6pt plus 1pt minus 1pt}
\setlength{\intextsep}{6pt plus 1pt minus 1pt}
\setlength{\abovecaptionskip}{3pt}
\setlength{\belowcaptionskip}{2pt}


\renewcommand{\thefigure}{S\arabic{figure}}
\renewcommand{\thetable}{S\arabic{table}}

\begin{document}
\bstctlcite{IEEEexample:BSTcontrol}

\title{Supplementary Material: A Bedside Decision Support System for Disorders of Consciousness: Multi-Scale Temporal Dynamics and Multicenter Validation}

\author{Hongjie~Chen$^{1}$, Jianming~Chen$^{1}$, Guanyu~Xiong$^{2}$, Fang~Li$^{2}$, and Fang~Duan$^{1,*}$%
\thanks{$^{1}$Hongjie Chen, Jianming Chen, and Fang Duan are with the College of Information Science and Engineering, Huaqiao University, Xiamen 361021, China.}%
\thanks{$^{2}$Guanyu Xiong and Fang Li are with the China Rehabilitation Research Center, Beijing Bo'ai Hospital, Capital Medical University, Beijing, China.}%
\thanks{$^{*}$Corresponding author: Fang Duan (e-mail: nkfetsh@gmail.com).}%
}

\markboth{IEEE JOURNAL OF BIOMEDICAL AND HEALTH INFORMATICS: SUPPLEMENTARY MATERIAL}%
{Chen et al.: Bedside Decision Support System for Disorders of Consciousness}

\maketitle

\begin{table*}[!ht]
\centering
\caption{De-Identified Individual Cohort Characteristics, Feature Profiles, and Cross-Center Benchmark Performance for the China Rehabilitation Research Center (CRRC / Beijing Bo'ai Hospital) Inpatient Cohort ($N=15$; 15 males; 6 MCS [including 1 EMCS], 9 VS/UWS; 11 ICH, 4 TBI; mean age $51.7 \pm 12.3$~years). Clean 6-channel records ($N=13$; 6 MCS [including 1 EMCS], 7 VS/UWS) exclude CRRC\_008 and CRRC\_009 due to continuous mechanical electrode drift on $O_2$.}
\label{tab:crrc_patient_cohort}
\vspace{-4pt}
\scriptsize
\renewcommand{\arraystretch}{0.92}
\setlength{\tabcolsep}{3.5pt}
\begin{tabular*}{\linewidth}{@{\extracolsep{\fill}}lccccccccc}
\hline
\textbf{Subject ID} & \textbf{Diagnosis} & \textbf{CRS-R} & \textbf{Lesion Laterality} & \textbf{Duration (s)} & \textbf{Retention} & $c_{\text{global}}$ & $-\mathrm{CV}_\theta$ & $\rho_1$ & $r_{\text{residual}}$ \\
\hline
CRRC\_001 & MCS & 11 & Bilateral intact & 605.0 & 98.7\% & 0.1834 & -1.3927 & +0.1118 & +0.5878 \\
CRRC\_002$^*$ & $\text{MCS}^+$$^*$ (EMCS) & 23 & Right hemisphere & 1236.0 & 98.7\% & 0.5343 & -0.3331 & +0.3201 & +0.2064 \\
CRRC\_003 & MCS & 10 & Left hemisphere & 1225.0 & 97.1\% & 0.8409 & -0.3682 & +0.1451 & +0.3646 \\
CRRC\_004 & MCS & 7 & Left hemisphere & 1221.0 & 100.0\% & 0.5993 & -0.3748 & -0.0378 & +0.1391 \\
CRRC\_005 & MCS & 17 & Bilateral intact & 845.0 & 100.0\% & 0.4906 & -0.6205 & +0.3850 & +0.3328 \\
CRRC\_006 & MCS & 9 & Right hemisphere & 1216.0 & 97.4\% & 0.5173 & -0.4568 & +0.5073 & +0.2188 \\
\hline
CRRC\_007 & VS/UWS & 7 & Right hemisphere & 1896.0 & 99.8\% & 0.8119 & -0.4731 & +0.0105 & +0.3776 \\
CRRC\_008$^\dagger$ & VS/UWS & 2 & Bilateral intact & 1213.0 & 100.0\% & 0.7253 & -0.6384 & +0.0337 & +0.3048 \\
CRRC\_009$^\dagger$ & VS/UWS & 3 & Bilateral intact & 1209.0 & 99.7\% & 0.6213 & -0.6091 & +0.0614 & +0.1660 \\
CRRC\_010$^\ddagger$ & VS/UWS & 2 & Bilateral intact & 637.0 & 21.4\% & 0.7294 & -0.4525 & +0.2401 & -0.0657 \\
CRRC\_011 & VS/UWS & 7 & Bilateral intact & 1237.0 & 99.0\% & 0.4154 & -0.4882 & +0.4003 & +0.3513 \\
CRRC\_012 & VS/UWS & 7 & Bilateral intact & 436.0 & 99.1\% & 0.5607 & -0.3266 & +0.0680 & +0.0339 \\
CRRC\_013 & VS/UWS & 7 & Left hemisphere & 2344.0 & 99.5\% & 0.5517 & -0.4399 & +0.1056 & +0.2702 \\
CRRC\_014 & VS/UWS & 7 & Right hemisphere & 1215.0 & 86.8\% & 0.7556 & -0.6032 & +0.6941 & +0.1504 \\
CRRC\_015 & VS/UWS & 4 & Left hemisphere & 1040.0 & 100.0\% & 0.6218 & -0.7362 & -0.0033 & +0.3001 \\
\hline
\multicolumn{10}{p{\dimexpr\linewidth-2\tabcolsep\relax}}{\footnotesize $^*$CRRC\_002 met Emergence from MCS (EMCS) criteria; retained in conscious class ($\text{MCS}^+$) for binary triage. $^\dagger$Mechanical $O_2$ electrode drift; excluded in Protocol 1 ($N=13$), retained in Protocol 2 ($N=15$). $^\ddagger$21.4\% retention due to motion; retained under intent-to-treat. Telemetry: NeuSen NSM2 / EEGCap32-pony (1000~Hz, 4-s epochs, CAR, F3, F4, C3, C4, O1, O2; IRB No.~2025-118-01~\cite{duan2026auditory}). Model comparisons across protocols appear in Table~\ref{tab:matched_comparison}; operating thresholds are detailed in Note~C.} \\
\hline
\end{tabular*}
\end{table*}

\begin{table*}[t]
\centering\footnotesize\vspace{-3pt}
\caption{Matched-input model comparisons (AUROC). All methods share feature columns, source-only parameter selection, scaling, direction coding, and source-training median thresholds within each protocol. Truncation is a fixed scoring-rule variant ($c=1.5$). Target standardization and direction alignment are unlabeled batch adaptations.}
\label{tab:matched_comparison}
\renewcommand{\arraystretch}{0.92}
\begin{tabular*}{\textwidth}{@{\extracolsep{\fill}}lccccc}
\hline
\textbf{Cohort and protocol} & \textbf{Linear} & \textbf{Clipped} & \textbf{Ridge} & \textbf{Logistic} & \textbf{LDA} \\
\hline
Mendeley, 25 outer folds & 0.9324 & 0.9291 & 0.8813 & 0.8964 & 0.8769 \\
CRRC 15, source frozen & 0.5741 & 0.5741 & 0.5556 & 0.5556 & 0.5556 \\
CRRC 15, target standardized & 0.5000 & 0.5000 & 0.5185 & 0.5185 & 0.5370 \\
CRRC 15, target aligned & 0.7778 & 0.7963 & 0.7407 & 0.7593 & 0.7593 \\
CRRC 13, target aligned & 0.7143 & 0.7143 & 0.7143 & 0.6905 & 0.6905 \\
Figshare 50, $D_1$, source frozen & 0.7691 & 0.7606 & 0.7759 & 0.7674 & 0.7742 \\
Figshare 50, $D_1$, target standardized & 0.7657 & 0.7606 & 0.7708 & 0.7589 & 0.7691 \\
\hline
\end{tabular*}
\end{table*}


\section*{Supplementary Clinical and Methodological Notes}

\subsection*{A. CRRC Cohort Recruitment and Ethics Compliance}
The prospective inpatient hospital cohort was evaluated at the China Rehabilitation Research Center (CRRC / Beijing Bo'ai Hospital, Capital Medical University Rehabilitation College). Inpatients presenting with disorders of consciousness secondary to severe acute brain injury (11 intracerebral hemorrhage [ICH], 4 traumatic brain injury [TBI]) were enrolled during neurorehabilitation. Ethical approval was granted by the Medical Ethics Committee of the China Rehabilitation Research Center (Approval No.~2025-118-01; cohort recruitment and auditory paradigm screening detailed in Duan et~al.~\cite{duan2026auditory}). Written informed surrogate consent was obtained from legally authorized representatives prior to telemetry recording in accordance with the Declaration of Helsinki. Diagnostic categorization into minimally conscious state (MCS/EMCS, $N=6$; 5 MCS and 1 EMCS) and vegetative state / unresponsive wakefulness syndrome (VS/UWS, $N=9$) was established through serial Coma Recovery Scale-Revised (CRS-R) assessments by certified neurorehabilitation clinicians (scores: 2--23). Recordings were conducted during stable post-acute neurorehabilitation in patients free from continuous acute ICU sedation. Telemetry was acquired from the cohort in~\cite{duan2026auditory}; daily medication logs were unavailable.

\subsection*{B. Telemetry Acquisition and Signal Preprocessing}
Bedside electroencephalography was acquired using a research-grade NeuSen NSM2 amplifier / EEGCap32-pony system configured at a native sampling rate of 1,000~Hz with active Ag/AgCl electrodes positioned according to the international 10--20 montage (19 channels; see Duan et~al.~\cite{duan2026auditory}). Continuous recordings spanned intermediate monitoring intervals (mean $19.53 \pm 7.96$~minutes; range 7.27 to 39.07~minutes). Raw continuous signals were re-referenced to the Common Average Reference (CAR) and segmented into non-overlapping 4-second blocks (4,000 samples per block). 

To prevent data leakage, artifact screening operated strictly at the individual recording level: blocks with peak-to-peak amplitude at or above $500\,\mu\mathrm{V}$ were flagged as artifactual gaps ($\mathbb{I}_{\text{adj}} = 0$). Homologous electrodes matching the discovery montage (F3, F4, C3, C4, O1, O2) were extracted. 

Two subjects (CRRC\_008 and CRRC\_009) exhibited persistent high-impedance drift isolated to the occipital channel $O_2$. To assess clinical robustness across sensor conditions, two complementary validation protocols were evaluated:
\begin{enumerate}
    \item \textbf{Protocol 1 (Strict Clean 6-Channel Protocol, $N=13$)}: Excluded the two subjects with localized $O_2$ artifacts, evaluating all records with fully intact 6-channel arrays.
    \item \textbf{Protocol 2 (Full Sensitivity Protocol, $N=15$)}: Retained all 15 enrolled subjects under an intent-to-screen sensitivity protocol regardless of single-channel electrode impedance.
\end{enumerate}

\subsection*{C. Matched Cross-Center Preprocessing and Operating Points}
All matched comparisons use the same feature matrix and source training records for every model. The primary matrix is $[c_{\text{global}},-\mathrm{CV}_\theta,\rho_1,r_{\text{residual}}]$. Source-only inner three-fold cross-validation selects ridge and logistic regularization; LDA uses analytical shrinkage. Full-source feature signs are $\mathbf{s}_{\text{source}}=[+1,+1,+1,-1]^\top$. Model fitting uses source-standardized, sign-coded coordinates.

Frozen-source transfer reuses source means, scales, and signs. Target-standardized transfer computes means and scales from the unlabeled target batch while retaining source signs. CRRC target-aligned transfer uses the same target-standardized matrix and $\mathbf{s}_{\text{target}}=[-1,+1,+1,+1]^\top$ for all models. In raw standardized coordinates, this corresponds to $\operatorname{diag}(-1,+1,+1,-1)$ applied consistently to the target inputs of source-trained supervised models. Target labels enter only the final evaluation. The direction rule is specified for the subacute clinical setting; its prospective deployment requires the rule and calibration batch to be fixed before evaluating later patients.

The primary four-feature score is the arithmetic mean; the fixed-truncation variant averages $\operatorname{clip}(s_kz_k,-1.5,1.5)$ under identical calibration (0.9291 AUROC discovery, 0.7963 full CRRC vs. 0.9324 and 0.7778 unclipped; Table~\ref{tab:matched_comparison}). Primary decisions apply source-training score medians (CRRC truncated: sensitivity 0.6667, specificity 0.7778, balanced accuracy 0.7222; Figshare source-frozen threshold: balanced accuracy 0.5645, sensitivity 1.0000, specificity 0.1290). All transforms, traces, and input hashes are archived in \url{results/matched_model_comparison_v1/}; paired uncertainty applies 10,000 subject-level Dirichlet draws.

\subsection*{D. Section S-CDSS: Bedside Streamlined Implementation \& Latency Profiling}
To explore implementation feasibility, the four-feature temporal dynamic consensus pipeline was implemented as a streaming prototype (\url{docs/bedside_cdss_demo.html}, animation in \url{docs/bedside_playback_demo.gif}). The architecture accepts consecutive non-overlapping 4-s telemetry epochs ($f_s = 256\,\mathrm{Hz}$), subjecting each block to artifact gating ($<500\,\mu\mathrm{V}$). Temporal continuity tracking ensures that artifact gaps break physical adjacency ($\mathbb{I}_{\text{adj}} = 0$); feature estimation is gated until $\ge 10$ physically adjacent clean pairs accumulate. Data with $<10$ adjacent pairs yield \texttt{INSUFFICIENT\_DATA}. Source-frozen standardized coordinates ($z_k = s_k(x_k - \mu_k)/\sigma_k$) use archived thresholds ($\theta_0 = 0.0451$ for the raw score and $\theta_{1.5} = 0.0725$ for the clipped score), yielding illustrative candidate prioritization rather than definitive clinical referrals.

Cached-power replay of two source-training records (5,400 epochs per subject; P25 [VS] and P36 [MCS]) showed numerical agreement with offline batch processing ($|\Delta_{\text{feature}}| \le 2.77 \times 10^{-13}$, scoring error $0.000$; Table~\ref{tab:bedside_latency}). Cached-power replay timings were: cached power ingestion mean $0.0098\,\mathrm{ms}$ (p95 $0.011\,\mathrm{ms}$); four-feature extraction mean $2.13\,\mathrm{ms}$ (p95 $3.11\,\mathrm{ms}$); consensus scoring mean $0.011\,\mathrm{ms}$; and total evaluated update mean $2.15\,\mathrm{ms}$ (p95 $3.13\,\mathrm{ms}$), well below the 4,000\,ms inter-epoch window. Raw EEG preprocessing and clinical decision latency were not measured; no single-core or thread restriction was imposed. Table values average the two subject-level summaries (27 evaluated updates each), including their individual p95 values. Complete trajectories and test records are archived in \url{results/bedside_playback_demo/}.

\begin{table}[b]
\caption{Cached-Power Replay Timings and Verification (5,400 Epochs per Record; Two Records)}
\label{tab:bedside_latency}
\centering
\scriptsize
\setlength{\tabcolsep}{2pt}
\begin{tabular}{lrrrr}
\toprule
Processing Phase & Calls/record & Mean (ms) & Mean p95 (ms) & Max (ms) \\
\midrule
Cached Power Ingest & 5,373 & 0.0098 & 0.0111 & 0.7090 \\
Feature Extraction  &    27 & 2.1323 & 3.1117 & 3.9057 \\
Consensus Scoring   &    27 & 0.0110 & 0.0123 & 0.0164 \\
Evaluated Step Total&    27 & 2.1542 & 3.1330 & 3.9261 \\
\midrule
\multicolumn{5}{l}{Offline Consistency: Max feature $|\Delta| \le 2.77\times 10^{-13}$; Score error: $0.000$.} \\
\bottomrule
\end{tabular}
\end{table}

\begingroup
\interlinepenalty=10000
\bibliographystyle{IEEEtran}
\bibliography{references}
\endgroup

\end{document}
